\chapter{Incidental Findings}

\section*{Patient focus}
This chapter explains the examination or clinical problem in accessible language. Patients should use it to prepare questions and understand the purpose, experience, benefits, and limitations of imaging; local instructions and the official report take priority.


\label{ch:ch37_incidental_findings}

\section{Definition}
An incidental finding is an unexpected observation unrelated to the original clinical question. Its importance depends on probability of disease, consequences of missing it, patient age and risk, and the burden of further evaluation.

\subsection{The Scale of the Problem}
With the increasing use of cross-sectional imaging, incidentalomas are now identified in up to 40\% of CT examinations and 20\% of MRI studies. The challenge is to distinguish the small proportion that represents serious disease from the large proportion that is benign, without subjecting every patient to unnecessary follow-up.

\begin{pearlbox}[Clinical Pearl: The 40\% rule]
Nearly half of all CT scans will show at least one incidental finding. Developing a systematic approach to incidentalomas is not optional --- it is essential for every radiologist and clinician.
\end{pearlbox}

\section{Most Common Incidentalomas}

\begin{table}[h]
\centering
\caption{Most common incidentalomas, their frequency, and recommended management.}
\label{tab:incidentalomas}
\begin{tabular}{p{2.5cm} p{2.5cm} p{3.5cm} p{5cm}}
\toprule
\textbf{Location} & \textbf{Frequency} & \textbf{Typical Appearance} & \textbf{Management} \\
\midrule
Adrenal & 4--5\% of abdominal CT & Well-defined, $<4$~cm, homogeneous, $<10$~HU on unenhanced CT & No follow-up if $<4$~cm, $<10$~HU, homogeneous; CT/MRI characterisation if indeterminate \\
Liver lesion & 5--15\% of abdominal CT & Variable attenuation, enhancement pattern & Characterise by enhancement pattern (hemangioma, FNH, cyst); MRI if indeterminate \\
Lung nodule & 25--50\% of chest CT & Solid, sub-solid, or ground-glass & Fleischner Society guidelines based on size, morphology, and risk factors \\
Pituitary & 10--20\% of brain MRI & $<10$~mm, non-enhancing, within the sella & Incidental microadenoma: no follow-up unless hormonal symptoms; macroadenoma ($>10$~mm) requires endocrine evaluation \\
Renal & 10--15\% of abdominal CT & Solid or cystic lesion, normal enhancement & Bosniak classification for cystic; solid lesions $<3$~cm may be monitored (active surveillance) \\
Thyroid & 15--25\% of neck CT & Well-defined nodule, variable echogenicity on US & US correlation; FNA if suspicious features (TI-RADS 4--5) or $>1.5$~cm \\
\bottomrule
\end{tabular}
\end{table}

\section{Management Guidelines}

\subsection{Fleischner Society Guidelines (Pulmonary Nodules)}
The Fleischner Society provides evidence-based recommendations for management of incidentally detected pulmonary nodules in adults:

\begin{itemize}
    \item \textbf{Solid nodule $<6$~mm}: No routine follow-up needed (regardless of risk).
    \item \textbf{Solid nodule 6--8~mm}: Follow-up CT at 6--12 months, then at 18--24 months if stable.
    \item \textbf{Solid nodule $>8$~mm}: Follow-up CT at 3 months, PET-CT, or tissue sampling.
    \item \textbf{Sub-solid nodule $<6$~mm}: No follow-up.
    \item \textbf{Sub-solid nodule $\geq6$~mm}: Follow-up CT at 3--6 months to confirm persistence, then annually for 5 years.
    \item \textbf{Multiple solid nodules $<6$~mm} in low-risk patients: No follow-up.
\end{itemize}

\subsection{ACR Incidental Findings Committee Recommendations}
The ACR has published specific guidance for the management of incidentalomas in multiple organs:
\begin{itemize}
    \item \textbf{Adrenal}: Unenhanced CT $<10$~HU = lipid-rich adenoma; no follow-up if $<4$~cm and homogeneous.
    \item \textbf{Thyroid}: Ultrasound correlation for all solid nodules; apply TI-RADS scoring.
    \item \textbf{Kidney}: Bosniak classification for cystic lesions; active surveillance for small solid masses in patients with limited life expectancy.
    \item \textbf{Liver}: Characterise by enhancement pattern; MRI with hepatocyte-specific contrast for indeterminate lesions.
\end{itemize}

\begin{pearlbox}[Clinical Pearl: The adrenal incidentaloma algorithm]
Step 1: Measure unenhanced attenuation ($<10$~HU = lipid-rich adenoma). Step 2: Assess size ($<4$~cm = low risk). Step 3: Check hormone profile if suspicious features. Step 4: Characterise with MRI or washout CT if indeterminate. Most adrenal incidentalomas require no further action.
\end{pearlbox}

\section{Reporting}
Describe the finding accurately, state its significance, and give a proportionate recommendation with timeframe and modality when appropriate. If no follow-up is recommended, say so when that prevents unnecessary testing.

\subsection{When to Dismiss Without Follow-Up}
An incidentaloma can be dismissed without follow-up when:
\begin{itemize}
    \item It meets all criteria for a confirmed benign entity (e.g.\ simple cyst, lipoma, calcified granuloma).
    \item It is stable on prior imaging for an adequate interval.
    \item It is a recognised normal variant (see Chapter~\ref{ch:ch31_normal_variants}).
    \item The patient has a very limited life expectancy where follow-up would not alter management.
\end{itemize}

\begin{pearlbox}[Clinical Pearl: Saying ``no follow-up'' is a recommendation]
Explicitly stating that no follow-up is needed is one of the most valuable things a radiologist can do. It prevents unnecessary scans, patient anxiety, and healthcare costs.
\end{pearlbox}

\section{Patient Communication}
Patients often experience anxiety when told that an ``incidental finding'' has been detected. Clinicians should:
\begin{itemize}
    \item Explain that incidental findings are very common and most are benign.
    \item Provide the specific recommendation and timeline.
    \item Avoid unnecessary alarm while maintaining honesty about uncertainty.
    \item Reassure the patient that the finding is being monitored appropriately.
\end{itemize}

\begin{pearlbox}[Impression:]
Incidental findings are ubiquitous in modern imaging. A proportionate, evidence-based approach using validated guidelines (Fleischner, ACR Incidentalomas, Bosniak, TI-RADS) prevents both under-investigation of significant disease and over-investigation of benign findings. Always state when no follow-up is recommended.
\end{pearlbox}

\clearpage
